\clearpage
\section{Related Work}\label{section::relatedwork}


In recent years, passenger detection and counting methods based on deep learning have become an active research area. Researchers have explored various network architectures and techniques from different perspectives to address counting problems in public transportation and crowd scenarios.
\subsection{Traditional Crowd Counting Methods}
Before deep learning-based methods became mainstream, several traditional crowd counting methods existed. These methods typically relied on handcrafted features and classical machine learning algorithms.

In earlier studies, traditional approaches employed image-based local features for passenger detection and counting in buses, such as Histogram of Oriented Gradients (HOG) \cite{30_khan_passenger_2020} and Hough transform \cite{73_Huofu}. However, these methods performed poorly in complex scenes and dynamic environments, primarily because manually designed features possess limited generalization capability, making it challenging to manage variations in appearance, occlusions, and changes in illumination.

Traditional crowd counting methods have also explored various feature extraction and classification techniques. Haq et al. developed a hybrid computer vision approach for real-time bus passenger flow calculation, utilizing HOG feature extraction combined with Support Vector Machine (SVM) classification for head detection \cite{19_haq_fast_2020}. This method demonstrates the effectiveness of traditional machine learning approaches in controlled environments. Hoang and Jo proposed a joint component model for pedestrian detection in crowded scenes, employing extended feature descriptors including Multiple scale blocks-based Histograms of Oriented Gradients (MHOG) and parallelogram-based Haar-like features (PHF) \cite{21_hoang_joint_2016}. These traditional feature descriptors were designed to handle partially occluded pedestrians in dense crowd scenarios. Additionally, Cheng et al. provided a comprehensive survey of traditional image segmentation techniques, including histogram thresholding, edge detection, and region-based methods, which form the foundation for many early crowd counting approaches \cite{10_cheng_color_2001}. These traditional methods, while limited in complex scenarios, established important groundwork for the development of more advanced deep learning-based solutions.

\subsection{Density Estimation-based Methods}
Density estimation-based methods have demonstrated greater robustness in high-density crowd counting scenarios.
For instance, Anusiya et al. utilized Gaussian filters to create ground-truth density maps for crowd number estimation \cite{71_Density_estimation_method}. Moreno Rendon et al. proposed a computer vision-based approach employing deep learning techniques to estimate the number of passengers at TransMilenio bus stations. They introduced the TransMilenio-Javeriana dataset, containing nearly 900,000 head annotations from buses and stations, and trained a deep learning architecture specifically tailored for crowd counting tasks, generating density maps focused on head regions in various scenes \cite{41_moreno_rendon_passenger_2023}. Bai et al. presented a crowd density detection approach based on crowd aggregation patterns and multi-column convolutional neural networks, aiming to integrate real-time surveillance videos captured from multiple viewpoints and granularities \cite{06_bai_crowd_2021}.

Recent advances in density-based crowd counting have shown significant improvements in accuracy and robustness. Zhang et al. proposed MCNN (Multi-column Convolutional Neural Network) for crowd counting, which uses multiple columns with different receptive fields to capture features at different scales \cite{zhang2016single}. CSRNet (Congested Scene Recognition Network) by Li et al. introduced dilated convolutions to maintain spatial resolution while expanding the receptive field, achieving state-of-the-art performance on dense crowd counting \cite{li2018csrnet}. Context-aware crowd counting by Liu et al. incorporated spatial context information to improve counting accuracy in complex scenes \cite{liu2019context}. These methods have been particularly effective in high-density scenarios where traditional detection-based approaches struggle due to severe occlusions.

Furthermore, recent developments in density estimation have focused on addressing scale variation and perspective distortion. Sam et al. proposed Switching CNN for crowd counting, which uses a switching mechanism to select appropriate regressors for different crowd densities \cite{sam2017switching}. Sindagi and Patel provided a comprehensive survey of CNN-based crowd counting methods, highlighting the evolution from traditional approaches to deep learning solutions \cite{sindagi2017survey}. These advances have significantly improved the robustness of density-based crowd counting systems in challenging real-world scenarios.

Additionally, recent research has explored the application of crowd counting in specific transportation scenarios. Abdou and Erradi conducted a comprehensive survey of machine learning approaches for crowd counting, identifying key challenges and future research directions \cite{03_abdou_crowd_2020}. Their work highlights the importance of crowd counting in resource allocation and emergency management, particularly in public transportation contexts.

\subsection{Crowd Counting Based on Object Detection}
Different from density map estimation methods, crowd counting methods based on object detection directly regard each individual in a crowd as a target to be detected. These methods generally employ existing object detection models (such as Faster R-CNN, SSD, YOLO, etc.) to locate people in images and obtain the total number of people by counting the detected targets. The progress of deep learning technology has significantly improved the accuracy of object detection in complex scenes, making crowd counting methods based on object detection an important research direction.

Recent developments in object detection-based crowd counting have focused on improving accuracy in challenging scenarios. Redmon et al. introduced YOLO (You Only Look Once) as a real-time object detection system that has been widely adopted for crowd counting applications \cite{redmon2016you}. Bochkovskiy et al. proposed YOLOv4, which achieved significant improvements in speed and accuracy for real-time detection tasks \cite{bochkovskiy2020yolov4}. Wang et al. developed a multi-scale feature fusion approach for pedestrian detection in crowded scenes, addressing the challenge of scale variation in crowd counting \cite{wang2019multiscale}. These advances have made object detection-based methods increasingly viable for real-time crowd counting applications.

For example, Mukti et al. studied an automatic passenger counting (APC) technology based on face detection, which uses the Viola-Jones method for face detection. In addition, this method is integrated with an embedded system based on the Internet of Things for processing and data transmission \cite{43_mukti_pemodelan_2021}. The team of Islam also proposed a passenger seat availability framework based on face detection, which detects the face of each passenger in a bus and sends the data of empty and occupied seats to the data center \cite{24_islam_framework_2020}.

The Pronello team developed a low-cost APC system using a Raspberry Pi with a camera and the YOLOV5 object detection algorithm, and conducted a field experiment in cooperation with the cities of Turin and ASTI in Italy. The experiment shows that the low-cost system was able to achieve an accuracy of 72.27\% and 74.59\% \cite{49_pronello_evaluating_2023}.

Tsiktsiris et al. proposed an end-to-end, rotation-aware detection framework specifically designed for accurately detecting passengers using angle-oriented bounding boxes, especially for top-view fisheye images. It locates people in images and obtains the total number of people by counting the detected targets \cite{55_tsiktsiris_improving_2024}.


\subsection{Other Related Methods}
Zhang et al. proposed a passenger counting method for buses based on the Single Shot MultiBox Detector (SSD) and Kalman filtering. In their method, the SSD model was adapted into a binary classifier trained specifically on bus passenger datasets. This model detects passengers in each frame, while a Kalman filter tracks detected positions across consecutive frames. Ultimately, passenger boarding and alighting statistics are generated based on these trajectories \cite{66_zhang_bus_2020}. Ahmed introduced a novel technique that transforms regions of interest containing pedestrians into standardized shapes, then applies the Rotated Histogram of Oriented Gradients (RHOG) algorithm and an Support Vector Machine (SVM) classifier based on machine learning. This method significantly improves pedestrian detection performance across different overhead scenarios \cite{05_ahmed_person_2019}.

Recent advances in crowd counting have also explored the integration of multiple sensing modalities beyond traditional computer vision approaches. Chen et al. proposed a multi-modal approach combining computer vision with WiFi sensing for more robust passenger counting in public transportation \cite{chen2018multimodal}. This hybrid approach addresses the limitations of vision-only systems in challenging lighting conditions and provides complementary information for improved accuracy. Additionally, the use of attention mechanisms in crowd counting has gained significant attention. Wang et al. introduced an attention-based crowd counting network that adaptively focuses on relevant image regions, particularly effective in scenarios with varying crowd densities \cite{wang2018attention}. These developments demonstrate the ongoing evolution of crowd counting techniques toward more robust and accurate solutions that integrate multiple data sources and advanced neural architectures.

The application of crowd counting in transportation systems has also benefited from advances in overhead view detection. This approach is particularly effective in cluttered environments where overhead views provide better coverage and visibility. The comprehensive review by Radovan et al. provides insights into the effectiveness, scalability, and practicality of different passenger counting solutions, with a focus on image processing and machine learning techniques \cite{02_noauthor_review_nodate}.

Apart from vision-based methods, other sensor technologies have also been explored for passenger counting applications. Haig et al. examined an approach that detects real-time bus occupancy using mobile phones. Their participatory sensing system utilizes various built-in smartphone sensors, such as accelerometers and global positioning systems (GPS), to indirectly infer bus occupancy and passenger numbers through sensor data analysis \cite{18_haig_real-time_2021}. The Shibata research team studied a crowd density estimation system using cellular communication radio-wave sensing combined with deep learning, specifically aiming to address privacy concerns \cite{52_shibata_people_2019}. Michele Nitti et al. investigated a method utilizing WiFi transmission to monitor mobile devices, enabling measurement of passenger load on buses \cite{70_noauthor_iabacus_nodate}.

Liu et al. proposed an automatic passenger boarding and alighting detection algorithm for subways based on deep learning and optical flow. The algorithm consists of a MetroNext network, which integrates two Multiple Scales Attention Convolution (MSAC) blocks, along with an optical flow algorithm. This combination rapidly and accurately detects subway cars and passengers, predicting passenger movement directions to determine boarding or alighting activities \cite{38_liu_automatic_2021}.

%  以上部分来自之前的studienproject的正文内容复制粘贴，需要保留
% ----------------------------------------------------------------

Overall, crowd counting methods are migrating from traditional feature-based and regression models toward deep learning-based end-to-end solutions \cite{13_dogra_survey_2024}. In high-density scenes with severe occlusion, density map-based methods exhibit notable advantages \cite{17_gupta_crowd_2024}. For example, in the extremely crowded MOT20 scenario, Counting MOT integrates crowd density maps to identify occluded or missed targets and correct false detections, thereby substantially improving the performance of MOT \cite{125}. However, a survey by Gang et al. points out that, while there is substantial work on analysis, counting, and classification in crowd tasks, there remain considerable research opportunities in crowd detection and crowd simulation \cite{01_noauthor_deep_2024}.

\paragraph{Applications} Crowd counting technology demonstrates significant potential across multiple domains:
\begin{itemize}
  \item \textbf{Intelligent surveillance and safety}: In intelligent surveillance systems, crowd counting can be used to detect abnormal events, such as rapid short-term crowd aggregation or anomalous flow behavior, which is critical for public safety in high-density areas such as railway platforms \cite{01_noauthor_deep_2024}. Such systems are essential for understanding dynamic crowd behavior, implementing crowd control, and analyzing individual behavior to ensure public safety and security, optimize public spaces, and manage events effectively \cite{15_ganga_object_2024}, including monitoring motion, speed, and flow direction and detecting anomalous events such as fights or running \cite{03_abdou_crowd_2020}.
  \item \textbf{Transportation systems}: In public transit, crowd counting is widely used to estimate onboard passenger numbers (e.g., buses) to optimize routes and schedules, improve convenience, reduce costs, and save energy \cite{33or34_li_lightweight_2023}. For example, real-time detection of crowded buses enables dynamic scheduling strategies on congested routes \cite{55_tsiktsiris_improving_2024}. In metro platforms, Transformer-based real-time passenger detection networks such as MPDNet \cite{64_yang_mpdnet_2022} have been proposed to obtain timely and accurate passenger flow data, supporting operational decisions and ensuring safe and efficient operations.
  \item \textbf{Public infrastructure planning}: Crowd counting data can be used to assess flow patterns in various public places, thereby assisting infrastructure design and expansion. By tracking ingress and egress at shopping malls, transit hubs, and airports, one can evaluate the efficiency of entrance and exit configurations. Optimizing these layouts improves flow efficiency and safety while minimizing wait times \cite{17_gupta_crowd_2024}.
  \item \textbf{In summary}, by providing precise information about crowd density, behavioral patterns, and transportation utilization, crowd counting can effectively assist urban planners, public transit operators, and infrastructure managers in making better-informed decisions, thereby optimizing existing infrastructure and guiding future work.
\end{itemize}

\bigskip

\subsection{Introduction to Multi-Object Tracking}
MOT is a fundamental and critical task in computer vision. Its core objective is to simultaneously localize multiple objects of interest (e.g., pedestrians, vehicles, animals) in video sequences and to assign and maintain a unique identity (ID) for each, thus constructing complete motion trajectories over time \cite{103,116,152}. Unlike single-object tracking (SOT), MOT must address target entries and exits, inter-object similarity, and complex interactions \cite{86}.

MOT is valuable across a wide array of real-world applications, including but not limited to: intelligent transportation and autonomous driving (real-time monitoring of road vehicles and pedestrians to enhance safety and efficiency \cite{162}); video surveillance and public safety (crowd counting, behavior analysis, and incident response \cite{60_wang_catrack_2025}); robot navigation and human–robot interaction (enabling machines to perceive and track dynamic entities in their environment \cite{160}); sports analytics (analyzing athlete trajectories to support tactics and performance evaluation \cite{200}); and wildlife management (assessing wildlife populations to facilitate ecological conservation \cite{169}).

\subsubsection{Analysis of Detector Components}

In MOT systems, detectors and trackers play different roles, but their performance is interdependent and jointly determines overall effectiveness.

\paragraph{Detector}
Before discussing specific MOT algorithms, the detector is a key component that identifies and precisely localizes all objects of interest in each frame, providing bounding boxes, categories, and confidence scores. Detector performance directly affects subsequent tracking accuracy, since false negatives can cause track loss and false positives can introduce spurious trajectories \cite{169}. Common detector choices and characteristics include:
\begin{itemize}
    \item Two-stage detectors: Generate region proposals followed by classification and bounding-box regression. Faster R-CNN is representative, renowned for high accuracy but relatively slow speed \cite{88}.
    \item One-stage detectors: Predict bounding boxes and categories directly from images without region proposals, typically offering significant speed advantages for real-time applications \cite{145}.
    \item YOLO family: As a leading one-stage detector family, YOLO (from YOLOv3 to YOLOv10 and even YOLOv12) strikes an excellent balance between speed and accuracy. For instance, YOLOv5 is widely used in vehicle recognition \cite{121} and crowd detection \cite{105}. YOLOv8 improves pedestrian detection accuracy and reduces model size via techniques such as Soft-NMS and GhostConv \cite{190}. YOLOv9 integrates optical flow to enhance tracking accuracy and robustness, especially under occlusion, appearance variation, and fast motion \cite{141}. YOLOX combines efficiency with strong detection performance and is adopted by leading trackers such as ByteTrack \cite{132}.
\end{itemize}

\subsubsection{Mainstream MOT Frameworks}

Currently, MOT algorithms can be broadly categorized into Tracking-by-Detection (TBD) and Joint Detection and Embedding (JDE). In recent years, with the advent of Transformers, several new end-to-end approaches have also emerged.

\paragraph{Tracking-by-Detection}
The TBD paradigm is among the earliest and most widely adopted. It decomposes MOT into two sequential stages: object detection and data association. A pre-trained detector first identifies and localizes targets in each frame independently, and an association algorithm then matches detections with existing tracks to maintain identity continuity \cite{152}.

Representative methods include: SORT, which pioneered the use of a Kalman filter for motion prediction and the Hungarian algorithm for IoU-based association to achieve simple, efficient real-time tracking \cite{169}; and DeepSORT, which augments SORT with deep appearance embeddings (Re-ID), combining motion (Kalman filter) and appearance (cosine similarity) with a cascaded matching strategy to significantly reduce ID switches and improve robustness and accuracy in complex scenes \cite{211}. Advantages include a clear, modular architecture that allows independent optimization of the detector and tracker; stronger detectors directly boost tracking performance. Limitations include the separation between detection and tracking, which can reduce efficiency and hinder true end-to-end optimization; under heavy occlusion or target disappearance, reliance on detections can lead to fragmented tracks and identity switches \cite{175}.

\paragraph{Joint Detection and Embedding}
JDE addresses TBD limitations by integrating object detection and Re-ID (appearance feature learning) into a unified network for joint training, aiming to improve efficiency and enhance tracking via shared feature learning \cite{173,132}.

Representative methods include: FairMOT, a CenterNet-based joint detection-and-tracking framework with two homogeneous branches for detection and Re-ID, which finds center-based embeddings more suitable for tracking than anchor-based embeddings \cite{208}; and ByteTrack, which leverages YOLOX as the detector and introduces a simple yet effective association strategy that fully exploits low-confidence detections, effectively handling occlusion and track fragmentation and achieving state-of-the-art results on MOTChallenge \cite{110}. Advantages include improved computational efficiency and runtime via shared backbones, better suitability for real-time deployment, and the potential for end-to-end optimization. Limitations include potential conflicts between detection and embedding objectives that may impair detection performance; in crowded scenes, anchor-based detectors can exhibit ambiguity when learning Re-ID features \cite{173}.

\paragraph{Emerging methods (Transformers, attention, Siamese networks, GNNs)}
With the evolution of deep learning architectures, emerging methods seek deeper integration of detection, tracking, and association to realize truly end-to-end tracking and reduce reliance on heuristic matching \cite{108}. Representative examples include: MOTR, which extends DETR by introducing track queries to model tracked instances across frames and achieve end-to-end tracking through iterative updates and predictions \cite{131}; TransTrack \cite{97} and TrackFormer \cite{164}, which leverage Transformer query–key mechanisms to decouple detection and association; attention-based VGT-MOT \cite{85}, which predicts target locations via inter-frame attention and adaptively fuses motion and appearance similarities to handle occlusion; Siamese-based methods such as SiamMOT \cite{151}; and graph neural networks (GNNs), which model complex spatiotemporal relations by treating detections or tracks as nodes and affinities as edges, e.g., GMTracker \cite{209}, which emphasizes inter-frame matching with intra-frame context. These methods promise more robust end-to-end training and inference and greater potential under complex, non-linear motion, but they are computationally demanding, require longer training, and depend on large datasets; smaller datasets such as MOT17 may be insufficient to fully train query-based trackers \cite{155}.

\subsection{Challenges in Complex Real-World Scenarios}

Despite significant progress, MOT still faces numerous challenges in high-density, complex-background, and real-time scenarios—particularly on railway platforms \cite{59_wang_deep_2024}.

(1) Occlusion: In crowded platform environments, pedestrians are densely packed and frequently occlude one another or are occluded by fixed structures (e.g., pillars, billboards, train carriages). This makes it difficult for detectors to capture complete bounding boxes or extract stable appearance features, leading to missed detections and identity loss \cite{169}. Traditional TBD methods degrade notably in such cases. DeepSORT mitigates this to some extent via appearance features and cascaded matching but may still fail under prolonged or extreme occlusions \cite{134}.

(2) Identity switches: When multiple targets have similar appearance (e.g., uniforms), or when targets reappear after complex motion, fast movement, or long occlusion, trackers may confuse identities or assign the same ID to different trajectories, causing identity switches \cite{87}. IDSW is a key robustness indicator. DeepSORT reduces IDSW substantially through appearance embeddings, yet it remains a core challenge in crowded scenes \cite{134}.

(3) Real-time–accuracy trade-off: In settings needing immediate response (e.g., platforms), MOT systems must be real-time. If passenger flow cannot be updated promptly, early warnings or dispatch decisions may be delayed. Improving accuracy often requires more complex models and higher computational costs, which conflict with real-time constraints \cite{203}. YOLO detectors are well known for favorable speed–accuracy trade-offs and are commonly used for real-time MOT; however, careful balancing is still needed at high frame rates.

(4) High-density crowds and complex backgrounds: Platform scenes typically exhibit high crowd density alongside dynamic, complex backgrounds (e.g., train arrivals/departures, signage illumination changes, passenger baggage), complicating target–background separation and introducing inter-object motion interactions that are hard to learn \cite{86}. 

(5) Small-object detection: When cameras are distant or targets occupy a small fraction of the image, they become small objects with limited pixels and weak features, making them prone to missed detection \cite{28_kaur_comprehensive_2023}.

(6) Illumination changes: Platform lighting can be complex and time-varying (e.g., strong sunlight, dim nighttime lighting, shadows cast by passing trains), which destabilizes appearance features and degrades detector and tracker performance \cite{169}.

(7) Generalization and annotation dependence: Deep learning methods depend heavily on large-scale, high-quality annotations. Since MOT models are often trained on specific datasets, the diversity of real-world conditions (camera viewpoints, resolutions, environments) demands strong generalization \cite{02_noauthor_review_nodate}.

\subsection{Proposed Approach: YOLOv11 + DeepSORT}

In light of the above challenges, especially on railway platforms, we propose combining YOLOv11 with DeepSORT. This choice reflects a careful analysis of existing algorithms and their complementary strengths in handling real-world complexity.

\subsubsection{YOLOv11: Strong Detection Capability}
Following the evolution of the YOLO family (e.g., YOLOv5, YOLOv8, YOLOv9, YOLOv10, YOLOv12), further gains in detection accuracy, speed, and compactness are expected. The one-stage, fast, and accurate nature of YOLO makes it particularly effective for real-time tasks. It produces high-confidence bounding boxes that are crucial as inputs to subsequent tracking modules. We anticipate YOLOv11 to provide:
\begin{itemize}
    \item High-accuracy detection: Accurate pedestrian detection on railway platforms, reducing false negatives and false positives.
    \item Real-time processing: Meeting the real-time requirements of high-frame-rate video streams, even under fast motion.
    \item Small-object adaptability: Enhanced detection of distant or partially occluded small targets, such as head-level targets in our case.
    \item Robustness: Stable outputs under illumination changes and complex backgrounds common to platforms.
\end{itemize}

\subsubsection{DeepSORT: Robust Tracking}
DeepSORT is a recognized TBD tracker noted for robust occlusion handling and significantly reduced identity switches. It combines a Kalman filter for motion prediction with deep appearance embeddings for identity association \cite{174}.
\begin{itemize}
    \item Occlusion handling: With deep Re-ID features, targets reappearing after long occlusions can be re-identified, compensating for transient detector misses; we additionally train Re-ID models on our custom dataset.
    \item Identity consistency: Appearance features enhance discrimination among similar targets, reducing ID switches and ensuring trajectory continuity and accuracy.
    \item Predictive capability: The Kalman filter predicts next-frame positions from historical motion, helping maintain tracks during brief occlusions or at lower frame rates.
\end{itemize}

\subsubsection{Complementarity and Synergy}

Our YOLOv11 + DeepSORT design exploits their complementarity in MOT:
\begin{itemize}
    \item Strong detection with strong tracking: YOLOv11 supplies high-quality, efficient detections as a foundation; DeepSORT’s association mechanism effectively handles challenges such as occlusion to ensure trajectory continuity and identity consistency.
    \item Tailored to railway platforms: In high-density, heavily occluded, and dynamically changing platform environments, YOLOv11 maximizes recall, while DeepSORT’s occlusion handling and ID management mitigate identity switches and track fragmentation. The synergy improves accuracy and robustness while preserving real-time performance.
\end{itemize}

By combining these components, we aim to build an efficient, accurate, and robust MOT system to meet the demands of detection, tracking, and counting in challenging environments such as railway platforms.

